\documentclass[11pt,a4paper,twoside]{report}

% --- packages -----------------------------------------------------------------
\usepackage[utf8]{inputenc}
\usepackage[T1]{fontenc}
\usepackage{lmodern}
\usepackage[a4paper,margin=2.4cm]{geometry}
\usepackage{microtype}
\usepackage{xcolor}
\usepackage{listings}
\usepackage{hyperref}
\usepackage{enumitem}
\usepackage{tabularx}
\usepackage{booktabs}
\usepackage{longtable}
\usepackage{titlesec}
\usepackage{fancyhdr}
\usepackage{parskip}
\usepackage{tikz}
\usetikzlibrary{arrows.meta,positioning,calc,shapes.geometric,fit,
                 decorations.pathreplacing,backgrounds,chains}
\usepackage{amsmath,amssymb,amsthm}
\usepackage{multicol}
\usepackage{float}
\usepackage{caption}
\usepackage{makeidx}

% --- colours (Tokyo Night palette) --------------------------------------------
\definecolor{tnbg}{HTML}{1A1B26}
\definecolor{tnfg}{HTML}{C0CAF5}
\definecolor{tnblue}{HTML}{7AA2F7}
\definecolor{tnred}{HTML}{F7768E}
\definecolor{tnyellow}{HTML}{E0AF68}
\definecolor{tngreen}{HTML}{73DACA}
\definecolor{tnpurple}{HTML}{BB9AF7}
\definecolor{tnorange}{HTML}{FF9E64}
\definecolor{tncomment}{HTML}{565F89}
\definecolor{linkcolor}{HTML}{2C5AA0}
\definecolor{lightgray}{HTML}{EAEAEA}
\definecolor{midgray}{HTML}{AAAAAA}

% --- hyperref -----------------------------------------------------------------
\hypersetup{
  colorlinks=true,
  linkcolor=linkcolor,
  urlcolor=linkcolor,
  citecolor=linkcolor,
  pdftitle={sotBSD System Manual},
  pdfauthor={Lucas Sotomayor (BASE4 Security)},
  pdfsubject={Comprehensive system manual for the sotBSD microkernel OS},
  pdfkeywords={microkernel, Rust, capability, exokernel, BSD, deception, sotFS}
}

% --- listings -----------------------------------------------------------------
\lstdefinelanguage{rustlang}{
  keywords={fn,let,mut,pub,struct,enum,impl,trait,if,else,for,while,loop,
            match,return,use,mod,as,where,self,Self,const,static,unsafe,
            extern,move,async,await,break,continue,crate,dyn,false,true,
            in,ref,type},
  keywordstyle=\color{tnblue}\bfseries,
  ndkeywords={u8,u16,u32,u64,usize,i8,i16,i32,i64,isize,bool,char,str,
              Option,Result,Vec,Box,String,CapId,Rights,SoHandle,Message},
  ndkeywordstyle=\color{tngreen}\bfseries,
  identifierstyle=\color{black},
  sensitive=true,
  comment=[l]{//},
  morecomment=[s]{/*}{*/},
  commentstyle=\color{tncomment}\itshape,
  stringstyle=\color{tnred},
  morestring=[b]",
  morestring=[b]',
}

\lstdefinelanguage{tlaplus}{
  keywords={MODULE,EXTENDS,CONSTANTS,VARIABLES,Init,Next,Spec,THEOREM,
            UNCHANGED,ENABLED,TRUE,FALSE,IF,THEN,ELSE,LET,IN,CHOOSE,
            DOMAIN,SUBSET,UNION,EXCEPT},
  keywordstyle=\color{tnblue}\bfseries,
  sensitive=true,
  comment=[l]{\\*},
  morecomment=[s]{(*}{*)},
  commentstyle=\color{tncomment}\itshape,
  stringstyle=\color{tnred},
  morestring=[b]",
}

\lstset{
  language=rustlang,
  basicstyle=\ttfamily\footnotesize,
  numbers=left,
  numberstyle=\tiny\color{tncomment},
  numbersep=8pt,
  frame=single,
  framesep=4pt,
  rulecolor=\color{lightgray},
  backgroundcolor=\color{lightgray!30},
  showstringspaces=false,
  breaklines=true,
  tabsize=2,
  captionpos=b,
  xleftmargin=1.5em,
  aboveskip=1em,
  belowskip=1em,
}

% --- titlesec -----------------------------------------------------------------
\titleformat{\chapter}[display]
  {\normalfont\huge\bfseries\color{tnblue}}
  {\chaptertitlename\ \thechapter}{20pt}{\Huge}
\titleformat{\section}
  {\normalfont\Large\bfseries\color{tnblue}}
  {\thesection}{1em}{}
\titleformat{\subsection}
  {\normalfont\large\bfseries\color{tnblue!80!black}}
  {\thesubsection}{1em}{}
\titleformat{\subsubsection}
  {\normalfont\normalsize\bfseries\color{tnblue!60!black}}
  {\thesubsubsection}{1em}{}

% --- header / footer ----------------------------------------------------------
\pagestyle{fancy}
\fancyhf{}
\fancyhead[LE]{\textit{sotBSD System Manual}}
\fancyhead[RO]{\textit{\leftmark}}
\fancyfoot[C]{\thepage}
\renewcommand{\headrulewidth}{0.3pt}

% --- theorem environments -----------------------------------------------------
\theoremstyle{definition}
\newtheorem{definition}{Definition}[chapter]
\newtheorem{theorem}{Theorem}[chapter]
\newtheorem{invariant}{Invariant}[chapter]

% --- convenience macros -------------------------------------------------------
\newcommand{\sotOS}{\texttt{sotOS}}
\newcommand{\sotBSD}{\texttt{sotBSD}}
\newcommand{\sotFS}{\texttt{sotFS}}
\newcommand{\LUCAS}{\textsc{lucas}}
\newcommand{\syscall}[1]{\texttt{#1}}
\newcommand{\reg}[1]{\texttt{#1}}
\newcommand{\hex}[1]{\texttt{0x#1}}
\newcommand{\file}[1]{\texttt{#1}}
\newcommand{\cmark}{\color{tngreen}\checkmark}
\newcommand{\TODO}[1]{\textcolor{tnred}{\textbf{TODO:} #1}}

% ==============================================================================
\begin{document}
% ==============================================================================

% --- Title page ---------------------------------------------------------------
\begin{titlepage}
\centering
\vspace*{3cm}

{\Huge\bfseries\color{tnblue} sotBSD System Manual\par}
\vspace{0.8cm}
{\Large\color{tnblue!70!black} The SOT Exokernel, sotBSD Microkernel,\\
 and sotFS Filesystem\par}
\vspace{0.5cm}
{\large Version 0.2.0 --- Codename: Project STYX\par}

\vspace{2cm}

\begin{tikzpicture}[
  layer/.style={draw=tnblue, fill=tnblue!8, minimum width=10cm,
                minimum height=0.9cm, font=\sffamily\small,
                rounded corners=3pt, thick},
]
  \node[layer] (hw) at (0,0) {Hardware (x86\_64, LAPIC, IOAPIC, PCI)};
  \node[layer] (kern) at (0,1.2) {SOT Microkernel (5 primitives, 12 SOT syscalls)};
  \node[layer, minimum width=4.8cm] (sotfs) at (-2.6,2.4) {sotFS};
  \node[layer, minimum width=4.8cm] (dec) at (2.6,2.4) {Deception Engine};
  \node[layer, minimum width=4.8cm] (lucas) at (-2.6,3.6) {LUCAS (Linux ABI)};
  \node[layer, minimum width=4.8cm] (rump) at (2.6,3.6) {rump / BSD};
  \node[layer] (app) at (0,4.8) {Applications (busybox, git, nano, links, wget)};
\end{tikzpicture}

\vspace{2cm}

{\large Lucas Sotomayor\\
\texttt{lucas@base4.io}\\
BASE4 Security\par}

\vspace{1cm}
{\normalsize 2026-04-10\par}
\end{titlepage}

% --- Front matter -------------------------------------------------------------
\pagenumbering{roman}
\tableofcontents
\newpage
\listoffigures
\newpage
\listoftables
\newpage
\lstlistoflistings
\newpage

% ==============================================================================
% PART I --- ARCHITECTURE
% ==============================================================================
\pagenumbering{arabic}
\part{Architecture}

% ------------------------------------------------------------------------------
\chapter{Introduction}
\label{ch:intro}
% ------------------------------------------------------------------------------

\section{What is sotBSD?}

\sotBSD{} is a verified microkernel operating system written from scratch in
Rust, designed around five core principles: minimal kernel surface, capability-based
security, transactional secure objects, deception-aware personality layers, and
formal verifiability. The system runs on the \texttt{x86\_64-unknown-none}
bare-metal target and boots via the Limine~v8 BIOS protocol.

The name reflects two layers: the SOT (Secure Object Transaction) exokernel
that provides the fundamental abstraction for all kernel objects, and the BSD
personality layer that presents a POSIX-compatible interface to userspace
programs. In practice, the system also hosts a full Linux ABI compatibility
layer called \LUCAS{} (Linux User-mode Compatibility And Syscall translation),
allowing real musl-static and glibc-linked binaries to run unmodified.

\section{Design Philosophy}

The architecture follows a strict separation of concerns:

\begin{enumerate}[leftmargin=1.4em]
  \item \textbf{Minimal kernel.} The microkernel implements exactly five
        primitives: scheduling, IPC, capabilities, frame allocation, and IRQ
        virtualization. Everything else---virtual memory, filesystems, drivers,
        networking---lives in userspace servers.
  \item \textbf{Capabilities as the sole access mechanism.} Every kernel object
        is referenced through a capability. There are no ambient authorities,
        no root user, no global namespaces.
  \item \textbf{Transactional objects.} The SOT exokernel wraps kernel objects
        in transactional envelopes with three consistency tiers, enabling
        atomic multi-object operations with provenance tracking.
  \item \textbf{Deception as a first-class primitive.} Capability interposition
        enables transparent rewriting of system call results, allowing the OS
        to present fabricated views of its own identity to adversaries.
  \item \textbf{Formal foundations.} Six TLA+ specifications cover the core
        subsystems, and the sotFS filesystem is verified through DPO rewriting
        rules with bounded treewidth guarantees.
\end{enumerate}

\section{Document Organization}

This manual is organized into five parts:

\begin{description}[leftmargin=1.4em]
  \item[Part~I: Architecture] (Chapters~\ref{ch:intro}--\ref{ch:memory})
    covers the system-level design, boot process, memory layout, and the
    SOT exokernel model.
  \item[Part~II: Subsystems] (Chapters~\ref{ch:caps}--\ref{ch:security})
    details the capability system, IPC mechanisms, scheduler, and security
    infrastructure.
  \item[Part~III: Personalities] (Chapters~\ref{ch:lucas}--\ref{ch:deception})
    describes the Linux ABI compatibility layer, the BSD/rump personality,
    and the deception engine.
  \item[Part~IV: Services \& Drivers] (Chapters~\ref{ch:sotfs}--\ref{ch:net})
    covers the sotFS filesystem, device drivers, and the networking stack.
  \item[Part~V: Operations] (Chapters~\ref{ch:build}--\ref{ch:formal})
    provides build instructions, debugging guides, a complete syscall
    reference, and formal verification details.
\end{description}

\section{Conventions Used in This Manual}

\begin{itemize}[leftmargin=1.4em]
  \item System call names are set in \syscall{monospace}.
  \item Register names use the \reg{rdi}, \reg{rsi}, \reg{rdx} convention.
  \item Hexadecimal addresses are prefixed: \hex{B00000}.
  \item Code listings use the Rust language unless otherwise noted.
  \item The symbol \cmark{} indicates a feature that has been verified
        end-to-end in the boot test pipeline.
\end{itemize}


% ------------------------------------------------------------------------------
\chapter{System Architecture}
\label{ch:arch}
% ------------------------------------------------------------------------------

\section{Architectural Overview}

\sotBSD{} is structured as a layered system with a minimal microkernel at its
core. Figure~\ref{fig:arch-stack} shows the complete architecture stack from
hardware to applications.

\begin{figure}[H]
\centering
\begin{tikzpicture}[
  layer/.style={draw=tnblue, thick, minimum height=1.1cm,
                font=\sffamily\small, rounded corners=2pt, align=center},
  hw/.style={layer, fill=tnred!10, minimum width=13cm},
  kern/.style={layer, fill=tnblue!10, minimum width=13cm},
  mid/.style={layer, fill=tngreen!10},
  upper/.style={layer, fill=tnyellow!10},
  app/.style={layer, fill=tnpurple!10, minimum width=13cm},
  arr/.style={-{Stealth[length=3mm]}, thick, tnblue!70},
]
  % Hardware
  \node[hw] (hw) at (0,0) {%
    \textbf{Hardware}\\[-2pt]
    {\footnotesize x86\_64 CPU, LAPIC, IOAPIC, PCI bus, NVMe, xHCI, virtio}};

  % Kernel
  \node[kern] (kern) at (0,1.6) {%
    \textbf{SOT Microkernel} (Ring~0, ${\sim}$10K LOC)\\[-2pt]
    {\footnotesize Sched\,|\,IPC\,|\,Caps\,|\,Frames\,|\,IRQ virt\,|\,%
     12 SOT syscalls\,|\,Provenance}};

  % Middle layer
  \node[mid, minimum width=3.1cm] (vmm) at (-5,3.3)
    {\textbf{VMM}\\[-2pt]{\footnotesize Paging}};
  \node[mid, minimum width=3.1cm] (sotfs) at (-1.7,3.3)
    {\textbf{sotFS}\\[-2pt]{\footnotesize VFS + ObjStore}};
  \node[mid, minimum width=3.1cm] (dec) at (1.7,3.3)
    {\textbf{Deception}\\[-2pt]{\footnotesize Interposition}};
  \node[mid, minimum width=3.1cm] (drv) at (5,3.3)
    {\textbf{Drivers}\\[-2pt]{\footnotesize blk/net/usb}};

  % Upper layer
  \node[upper, minimum width=4.2cm] (lucas) at (-3.3,5)
    {\textbf{LUCAS}\\[-2pt]{\footnotesize Linux ABI (127+ syscalls)}};
  \node[upper, minimum width=4.2cm] (rump) at (1.1,5)
    {\textbf{rump kernel}\\[-2pt]{\footnotesize NetBSD personality}};
  \node[upper, minimum width=3.6cm] (illumos) at (5,5)
    {\textbf{Illumos}\\[-2pt]{\footnotesize SMF/Crossbow/FMA}};

  % Applications
  \node[app] (apps) at (0,6.6) {%
    \textbf{Applications}\\[-2pt]
    {\footnotesize busybox, git, nano, links, wget, alpine packages,
     custom services}};

  % Arrows
  \draw[arr] (hw) -- (kern);
  \draw[arr] (kern) -- (vmm);
  \draw[arr] (kern) -- (sotfs);
  \draw[arr] (kern) -- (dec);
  \draw[arr] (kern) -- (drv);
  \draw[arr] (vmm) -- (lucas);
  \draw[arr] (sotfs) -- (lucas);
  \draw[arr] (sotfs) -- (rump);
  \draw[arr] (dec) -- (lucas);
  \draw[arr] (dec) -- (rump);
  \draw[arr] (drv) -- (illumos);
  \draw[arr] (lucas) -- (apps);
  \draw[arr] (rump) -- (apps);
  \draw[arr] (illumos) -- (apps);
\end{tikzpicture}
\caption{sotBSD architecture stack. The microkernel provides five primitives;
         everything above runs in Ring~3 as capability-holding services.}
\label{fig:arch-stack}
\end{figure}

\section{The Five Kernel Primitives}

The microkernel implements exactly five primitives. All other OS functionality
is built on top of these through userspace services communicating via IPC.

\begin{description}[leftmargin=1.4em]
  \item[Scheduling.] A four-level priority Multi-Level Queue (MLQ) scheduler
    with per-core Chase-Lev work-stealing deques. Priority levels are:
    \textit{realtime} (0), \textit{high} (1), \textit{normal} (2), and
    \textit{low} (3). The scheduler supports up to 16 CPUs via SMP with
    LAPIC timer-driven preemption and IPI-based cross-core wakeup.

  \item[IPC.] Two mechanisms: synchronous endpoints (L4-style rendezvous
    where sender blocks until receiver calls) and asynchronous channels
    (lock-free SPSC ring buffers for high-throughput data transfer). IPC
    also supports capability transfer in the message registers.

  \item[Capabilities.] All kernel objects are accessed through capabilities.
    A Capability Derivation Tree (CDT) tracks parent-child relationships
    for delegation and epoch-based O(1) revocation. A rights bitmask
    (Read/Write/eXecute/Grant/Revoke) provides monotonic attenuation.

  \item[Frame Allocation.] A bitmap-based physical frame allocator manages
    all physical memory. A slab allocator built on top provides efficient
    allocation for fixed-size kernel objects (thread control blocks,
    capability entries, endpoint structures).

  \item[IRQ Virtualization.] Hardware interrupts are demultiplexed by the
    kernel and delivered as IPC messages to registered userspace handlers.
    This allows device drivers to run entirely in Ring~3 without any
    kernel-mode driver code.
\end{description}

\section{The SOT Exokernel Layer}

On top of the five primitives, the SOT (Secure Object Transaction) layer
provides 12 additional syscalls (numbers 300--311) that implement:

\begin{itemize}[leftmargin=1.4em]
  \item \textbf{Secure Objects (SO):} opaque, typed blobs that serve as the
    universal container for all named kernel resources.
  \item \textbf{Transactions:} three tiers of consistency---Tier~0 (RCU
    read-only snapshots), Tier~1 (WAL single-object writes), and Tier~2
    (two-phase commit for multi-object atomic operations).
  \item \textbf{Provenance:} per-CPU SPSC ring buffers (4096 entries each)
    that record a SHA-256 hash-linked DAG of all SO mutations, enabling
    full audit trails and tamper detection.
\end{itemize}

\subsection{Category-Theoretic Foundations}

The SOT model is grounded in five categories with three functors:

\begin{definition}[SOT Categories]
The system defines five categories:
\begin{align}
  \mathbf{SO}   &: \text{Secure Objects (morphisms = capability derivations)} \\
  \mathbf{Dom}  &: \text{Protection Domains (morphisms = domain transitions)} \\
  \mathbf{Chan} &: \text{Channels (morphisms = message flows)} \\
  \mathbf{Tx}   &: \text{Transactions (morphisms = commit/abort)} \\
  \mathbf{Prov} &: \text{Provenance (morphisms = hash-chain links)}
\end{align}
\end{definition}

Three functors connect these categories:
\begin{itemize}[leftmargin=1.4em]
  \item $F_{\text{Access}} : \mathbf{Dom} \to \mathbf{SO}$ maps domain
    transitions to capability derivations.
  \item $F_{\text{Prov}} : \mathbf{Tx} \to \mathbf{Prov}$ maps transaction
    commits to provenance entries.
  \item $F_{\text{Deception}} : \mathbf{SO} \to \mathbf{SO}$ is an
    endofunctor that transparently rewrites SO contents for deception.
\end{itemize}

\section{The 12 SOT Syscalls}

Table~\ref{tab:sot-syscalls} lists the complete SOT syscall interface.

\begin{table}[H]
\centering
\caption{SOT exokernel syscalls (numbers 300--311).}
\label{tab:sot-syscalls}
\begin{tabularx}{\textwidth}{c l l X}
\toprule
\textbf{\#} & \textbf{Name} & \textbf{Args} & \textbf{Description} \\
\midrule
300 & \syscall{so\_create}    & type, size, data  & Create a new Secure Object \\
301 & \syscall{so\_read}      & so\_cap, off, len & Read bytes from an SO \\
302 & \syscall{so\_write}     & so\_cap, off, data & Write bytes to an SO \\
303 & \syscall{so\_destroy}   & so\_cap           & Destroy an SO and free storage \\
304 & \syscall{cap\_derive}   & src\_cap, rights   & Derive child cap with restricted rights \\
305 & \syscall{cap\_revoke}   & cap               & Revoke cap and all descendants in CDT \\
306 & \syscall{chan\_send}     & chan\_cap, msg     & Send message on async channel \\
307 & \syscall{chan\_recv}     & chan\_cap          & Receive message from async channel \\
308 & \syscall{domain\_enter}  & dom\_cap          & Enter a protection domain \\
309 & \syscall{domain\_yield}  &                   & Yield back to calling domain \\
310 & \syscall{tx\_commit}     & tx\_handle        & Commit a transaction (Tier 1 or 2) \\
311 & \syscall{tx\_abort}      & tx\_handle        & Abort a transaction, discard writes \\
\bottomrule
\end{tabularx}
\end{table}

% ------------------------------------------------------------------------------
\chapter{Boot Process}
\label{ch:boot}
% ------------------------------------------------------------------------------

\section{Boot Sequence Overview}

The system boots through four distinct phases, illustrated in
Figure~\ref{fig:boot-sequence}.

\begin{figure}[H]
\centering
\begin{tikzpicture}[
  phase/.style={draw=tnblue, thick, fill=tnblue!8, minimum width=3cm,
                minimum height=1.4cm, rounded corners=3pt,
                font=\sffamily\small, align=center},
  arr/.style={-{Stealth[length=3mm]}, very thick, tnblue!70},
  note/.style={font=\sffamily\scriptsize\color{tncomment}, align=left},
]
  \node[phase] (bios) at (0,0)   {\textbf{BIOS/Limine}\\firmware};
  \node[phase] (kern) at (4,0)   {\textbf{Kernel}\\kmain()};
  \node[phase] (init) at (8,0)   {\textbf{Init}\\services/init};
  \node[phase] (svc) at (12,0)   {\textbf{Services}\\VMM, FS, Net};
  \node[phase] (shell) at (12,-2.5) {\textbf{LUCAS Shell}\\interactive};

  \draw[arr] (bios) -- (kern);
  \draw[arr] (kern) -- (init);
  \draw[arr] (init) -- (svc);
  \draw[arr] (svc) -- (shell);

  \node[note, above=0.3cm of bios] {Limine v8 BIOS\\loads ELF kernel};
  \node[note, above=0.3cm of kern] {GDT, IDT, serial\\frame alloc, SMP};
  \node[note, above=0.3cm of init] {spawn VMM, kbd\\net, NVMe, xHCI};
  \node[note, above=0.3cm of svc] {register services\\via SvcRegister};
  \node[note, right=0.3cm of shell] {busybox, git,\\nano, wget};
\end{tikzpicture}
\caption{Boot sequence from firmware through kernel initialization to the
         interactive \LUCAS{} shell.}
\label{fig:boot-sequence}
\end{figure}

\subsection{Phase 1: Firmware and Bootloader}

The Limine~v8 BIOS bootloader loads the kernel ELF binary from a GPT+FAT32
disk image. The kernel must be compiled as \texttt{ET\_EXEC} (not
\texttt{ET\_DYN}), which requires the \texttt{-C relocation-model=static}
flag. Limine passes control to \texttt{kmain()} with a set of boot
information structures describing available memory, framebuffer, and the
CPIO initrd archive.

The boot configuration lives in \file{boot/limine.conf}:
\begin{lstlisting}[language={},caption={Limine boot configuration (excerpt).}]
serial: yes
timeout: 0

/sotOS
  protocol: limine
  kernel_path: boot():/sotos-kernel
  module_path: boot():/initrd.cpio
\end{lstlisting}

\subsection{Phase 2: Kernel Initialization}

The kernel initialization sequence in \texttt{kmain()} follows a strict order:

\begin{enumerate}[leftmargin=1.4em]
  \item \textbf{Serial FIRST.} Initialize COM1 serial output before any
    assertions or debug prints. This is a critical invariant---panics
    before serial is initialized produce no diagnostics.
  \item \textbf{GDT and IDT.} Load the Global Descriptor Table and
    Interrupt Descriptor Table. Configure IST entries: \#DF (index~0),
    \#GP (index~2), misc (index~3). \textit{Note:} \#PF does not use IST
    to avoid corruption when multiple threads fault simultaneously.
  \item \textbf{Frame allocator.} Initialize the bitmap-based physical
    frame allocator from the Limine memory map.
  \item \textbf{Slab allocator.} Initialize the kernel slab allocator on
    top of the frame allocator.
  \item \textbf{Capability tables.} Initialize the global capability table
    and CDT structures.
  \item \textbf{Scheduler.} Initialize per-CPU run queues and the
    Chase-Lev work-stealing deques.
  \item \textbf{SMP.} Wake Application Processors (APs) via INIT-SIPI-SIPI.
    Enable SSE on all CPUs (CR0.EM=0, CR4.OSFXSR, CR4.OSXMMEXCPT).
  \item \textbf{LAPIC timer.} Calibrate and start the Local APIC timer for
    preemptive scheduling.
  \item \textbf{Init process.} Load the \texttt{init} service from the
    initrd CPIO archive and create the first userspace thread.
\end{enumerate}

\subsection{Phase 3: Init Service}

The \file{services/init} process is the first userspace program. It reads
BootInfo from \hex{B00000} to obtain root capabilities, then spawns the
core services:

\begin{itemize}[leftmargin=1.4em]
  \item \textbf{VMM server:} manages page tables and handles \#PF
    (page faults) for all user processes.
  \item \textbf{Keyboard driver:} reads PS/2 scan codes and distributes
    input events via a shared ring buffer at \hex{510000}.
  \item \textbf{Network service:} implements the IP/TCP/UDP/ICMP stack.
  \item \textbf{NVMe driver:} communicates with NVMe SSDs via PRP lists
    and IPC.
  \item \textbf{xHCI driver:} manages USB host controllers and HID
    devices.
\end{itemize}

Each service registers itself with the kernel's service registry via
\syscall{SvcRegister} (syscall~130), making its endpoint discoverable by
other services through \syscall{SvcLookup} (syscall~131).

\subsection{Phase 4: Shell and Applications}

After services are running, init launches the \LUCAS{} shell, which provides
a POSIX-like command-line interface supporting pipes (up to 10 commands),
background jobs, environment variables, and scripting. Real Linux binaries
such as \texttt{busybox}, \texttt{git}, \texttt{nano}, and \texttt{links}
run through the Linux ABI compatibility layer.


% ------------------------------------------------------------------------------
\chapter{Memory Architecture}
\label{ch:memory}
% ------------------------------------------------------------------------------

\section{Physical Memory Management}

Physical memory is managed by a bitmap-based frame allocator. Each bit in the
bitmap represents a single 4\,KiB frame. The allocator is initialized from
the Limine memory map, which marks usable regions. A spinlock
(\texttt{spin::Mutex}) protects concurrent access.

The slab allocator provides efficient allocation for fixed-size kernel objects,
reducing fragmentation for frequently allocated structures such as thread
control blocks and capability entries.

\section{Virtual Address Space Layout}

Each process has its own address space (CR3). The init process uses the
layout shown in Table~\ref{tab:addr-layout}.

\begin{table}[H]
\centering
\caption{Init process virtual address space layout.}
\label{tab:addr-layout}
\begin{tabularx}{\textwidth}{l l X}
\toprule
\textbf{Address} & \textbf{Size} & \textbf{Content} \\
\midrule
\hex{200000}   & variable & ELF \texttt{.text} (kernel or service code) \\
\hex{510000}   & 4\,KiB  & Keyboard ring buffer \\
\hex{520000}   & 4\,KiB  & Mouse ring buffer \\
\hex{900000}   & 16\,KiB & Stack (4 pages) \\
\hex{A00000}   & 4\,KiB  & SPSC ring (IPC) \\
\hex{B00000}   & 4\,KiB  & BootInfo page (root caps) \\
\hex{B80000}   & 4\,KiB  & vDSO page (R+X) \\
\hex{C00000}+  & variable & Virtio/MMIO device regions (UC mapped) \\
\hex{D00000}+  & variable & ObjectStore + VFS data \\
\hex{E30000}+  & variable & Child process stacks \\
\hex{2000000}  & 1\,MiB  & \texttt{brk} heap \\
\hex{3000000}  & variable & \texttt{mmap} region \\
\hex{6000000}  & variable & \texttt{INTERP\_LOAD\_BASE} (dynamic linker) \\
\bottomrule
\end{tabularx}
\end{table}

\section{Copy-on-Write Fork}

The \syscall{SYS\_AS\_CLONE} (syscall~125) implements real copy-on-write fork:

\begin{enumerate}[leftmargin=1.4em]
  \item The \texttt{clone\_cow()} function walks all user-space PTEs in the
    parent's page tables.
  \item Both parent and child PTEs are marked read-only.
  \item A per-frame reference count (\texttt{FRAME\_REFCOUNT}) is incremented.
  \item On write fault (error code \texttt{\& 0x03 == 0x03}), the VMM
    allocates a new frame, copies the data via HHDM (Higher Half Direct
    Map), and remaps the page as writable.
\end{enumerate}

A critical implementation detail: \#PF (page fault) must \emph{not} use an
IST (Interrupt Stack Table) entry, because shared IST stacks corrupt when
multiple threads fault simultaneously during CoW fork. Instead, \#PF uses
the per-thread kernel stack (TSS RSP0).

\section{Guard Pages and W\^{}X}

Every thread stack is bounded by guard pages---unmapped pages that trigger
immediate \#PF on stack overflow rather than silently corrupting adjacent
memory. The kernel enforces W\^{}X (Write XOR Execute): no page may be
simultaneously writable and executable. The \syscall{SYS\_PROTECT}
(syscall~134) provides \texttt{mprotect}-like functionality with W\^{}X
enforcement.


% ==============================================================================
% PART II --- SUBSYSTEMS
% ==============================================================================
\part{Subsystems}

% ------------------------------------------------------------------------------
\chapter{Capability System}
\label{ch:caps}
% ------------------------------------------------------------------------------

\section{Capability Model}

Every kernel object is accessed through a capability. A capability is a
protected reference that bundles an object identity with an access rights
mask. There are no ambient authorities in the system---the only way to
perform any operation is to present a valid capability.

\subsection{CapId Format}

A \texttt{CapId} is a 32-bit value composed of two fields:

\begin{center}
\begin{tikzpicture}[font=\ttfamily\small]
  \draw[thick] (0,0) rectangle (8,0.8);
  \draw[thick] (5,0) -- (5,0.8);
  \node at (2.5,0.4) {index (20 bits)};
  \node at (6.5,0.4) {generation (12 bits)};
  \node[above] at (0,0.8) {\footnotesize bit 31};
  \node[above] at (5,0.8) {\footnotesize bit 12};
  \node[above] at (8,0.8) {\footnotesize bit 0};
\end{tikzpicture}
\end{center}

The 20-bit index allows up to $2^{20} = 1{,}048{,}576$ simultaneous
capabilities. The 12-bit generation counter prevents use-after-free
attacks: when a capability is revoked, the generation is bumped, and any
stale \texttt{CapId} with the old generation is rejected.

\subsection{Rights Bitmask}

Each capability carries a 5-bit rights mask:

\begin{table}[H]
\centering
\begin{tabular}{c l l}
\toprule
\textbf{Bit} & \textbf{Name} & \textbf{Meaning} \\
\midrule
0 & Read    & Read object contents \\
1 & Write   & Modify object contents \\
2 & Execute & Execute/invoke the object \\
3 & Grant   & Derive child capabilities \\
4 & Revoke  & Revoke child capabilities \\
\bottomrule
\end{tabular}
\end{table}

Rights are \emph{monotonically attenuated}: a derived capability can only
have a subset of its parent's rights. The \texttt{Rights::restrict()} method
enforces this invariant.

\section{Capability Derivation Tree}

The CDT (Capability Derivation Tree) tracks the parent-child relationship
between capabilities. Each capability entry stores a parent pointer. When
a capability is revoked via \syscall{cap\_revoke}, the revocation propagates
to all descendants through an epoch-based mechanism that achieves O(1)
amortized cost.

Figure~\ref{fig:cdt} shows a typical derivation tree with progressive
attenuation.

\begin{figure}[H]
\centering
\begin{tikzpicture}[
  capnode/.style={draw=tnblue, thick, fill=tnblue!10, rounded corners=3pt,
              minimum width=2.2cm, minimum height=0.8cm, align=center,
              font=\sffamily\small},
  arr/.style={-{Stealth[length=2.5mm]}, thick, tnblue!70},
  revoked/.style={capnode, draw=tnred, fill=tnred!10, dashed},
]
  % Root
  \node[capnode] (root) at (0,0) {Root\\{\scriptsize RWXGR}};

  % Level 1
  \node[capnode] (fs) at (-3,-1.8) {FS Server\\{\scriptsize RWXG-}};
  \node[capnode] (net) at (0,-1.8) {Net Server\\{\scriptsize RW-G-}};
  \node[capnode] (drv) at (3,-1.8) {Driver\\{\scriptsize R-X--}};

  % Level 2
  \node[capnode] (app1) at (-4.5,-3.6) {App~A\\{\scriptsize R-X--}};
  \node[capnode] (app2) at (-1.5,-3.6) {App~B\\{\scriptsize R----}};
  \node[revoked] (rev) at (1.5,-3.6) {Revoked\\{\scriptsize -----}};

  % Edges
  \draw[arr] (root) -- (fs) node[midway,left,font=\scriptsize] {derive};
  \draw[arr] (root) -- (net) node[midway,left,font=\scriptsize] {derive};
  \draw[arr] (root) -- (drv) node[midway,right,font=\scriptsize] {derive};
  \draw[arr] (fs) -- (app1) node[midway,left,font=\scriptsize] {restrict};
  \draw[arr] (fs) -- (app2) node[midway,right,font=\scriptsize] {restrict};
  \draw[arr] (net) -- (rev) node[midway,right,font=\scriptsize] {revoke};

  % Annotation
  \draw[decorate,decoration={brace,amplitude=8pt,raise=4pt},tncomment]
    (-5,-3.9) -- (-0.5,-3.9) node[midway,below=12pt,font=\sffamily\scriptsize\color{tncomment}]
    {Monotonic attenuation: children $\subseteq$ parent rights};
\end{tikzpicture}
\caption{Capability Derivation Tree with monotonic attenuation. The dashed
         node shows a revoked capability---all descendants are also invalidated.}
\label{fig:cdt}
\end{figure}

\section{Interposition Policies}

The capability system supports four interposition policies that control how
operations on a capability are handled:

\begin{description}[leftmargin=1.4em]
  \item[Passthrough.] The operation proceeds directly to the underlying
    object with no modification.
  \item[Inspect.] The operation is logged to the provenance subsystem
    before proceeding.
  \item[Redirect.] The operation is transparently forwarded to a
    different object (used by the deception engine).
  \item[Fabricate.] The result is synthesized without touching the real
    object (used for honeypot responses).
\end{description}

\section{Capability Derive and Revoke Example}

Listing~\ref{lst:cap-derive} shows how a service derives a restricted
capability and later revokes it.

\begin{lstlisting}[caption={Capability derivation and revocation.},
                   label={lst:cap-derive}]
use sotos_common::sys;
use sotos_common::cap::{CapId, Rights};

// Derive a read-only capability from our full-access cap
let full_cap: CapId = my_resource_cap;
let ro_rights = Rights::READ;
let ro_cap = sys::cap_derive(full_cap, ro_rights)
    .expect("cap_derive failed");

// Grant the read-only cap to a client via IPC
let mut msg = Message::new();
msg.set_cap(0, ro_cap);
sys::send(client_ep, &msg);

// Later: revoke the client's capability (and all its children)
sys::cap_revoke(ro_cap)
    .expect("cap_revoke failed");
// ro_cap and any caps derived from it are now invalid
\end{lstlisting}


% ------------------------------------------------------------------------------
\chapter{IPC Mechanisms}
\label{ch:ipc}
% ------------------------------------------------------------------------------

\section{Synchronous Endpoints}

Synchronous IPC follows the L4 rendezvous model: the sender blocks until
a receiver is ready, and vice versa. This provides zero-copy message
transfer with a context switch directly from sender to receiver.

\begin{figure}[H]
\centering
\begin{tikzpicture}[
  proc/.style={draw=tnblue, thick, fill=tnblue!10, minimum width=2.5cm,
               minimum height=3.5cm, rounded corners=3pt},
  state/.style={font=\sffamily\scriptsize, fill=white, draw=tnblue!50,
                rounded corners=2pt, minimum width=1.8cm,
                minimum height=0.5cm},
  arr/.style={-{Stealth[length=2.5mm]}, thick},
  msg/.style={font=\sffamily\scriptsize},
]
  % Sender
  \node[proc, label={above:\textbf{Sender}}] (S) at (0,0) {};
  \node[state] (s1) at (0,1) {Running};
  \node[state, fill=tnyellow!20] (s2) at (0,0) {Blocked};
  \node[state] (s3) at (0,-1) {Running};

  % Receiver
  \node[proc, label={above:\textbf{Receiver}}] (R) at (6,0) {};
  \node[state, fill=tnyellow!20] (r1) at (6,1) {Blocked};
  \node[state] (r2) at (6,0) {Running};
  \node[state, fill=tnyellow!20] (r3) at (6,-1) {Blocked};

  % Kernel
  \node[font=\sffamily\small\color{tncomment}] at (3,2) {\textbf{Kernel (rendezvous)}};

  % Arrows
  \draw[arr, tnblue] (s1.east) -- (r1.west)
    node[midway, above, msg] {send(ep, msg)};
  \draw[arr, tngreen] (r1.south east) -- (r2.north east)
    node[right, msg] {recv(ep)};
  \draw[arr, tnorange, dashed] (3,0.3) -- (3,-0.3)
    node[right, msg] {copy regs};
  \draw[arr, tnred] (r2.west) -- (s3.east)
    node[midway, below, msg] {reply(msg)};

  % Timeline
  \draw[very thin, tncomment] (-2,1.5) -- (-2,-1.5)
    node[below, font=\scriptsize] {time};
  \draw[-{Stealth[length=2mm]}, thin, tncomment] (-2,1.5) -- (-2,-1.5);
\end{tikzpicture}
\caption{Synchronous IPC flow. The sender blocks on \texttt{send()}, the
         receiver blocks on \texttt{recv()}, and the kernel performs a
         direct register-to-register copy at the rendezvous point.}
\label{fig:ipc-sync}
\end{figure}

A message consists of up to 8 general-purpose registers (64 bytes total of
inline data) plus optional capability transfers in dedicated capability
registers. For bulk data, the channel mechanism (Section~\ref{sec:async})
is used instead.

\subsection{IPC Timeout}

The \syscall{sys::call\_timeout(cap, msg, ticks)} interface allows a caller
to set a deadline on an IPC call. If the callee does not reply within the
specified number of scheduler ticks, the kernel wakes the blocked caller
with an error code.

\section{Asynchronous Channels}
\label{sec:async}

Asynchronous channels use lock-free single-producer single-consumer (SPSC)
ring buffers for high-throughput, non-blocking data transfer. Each channel
is a page-sized ring buffer (\hex{A00000}) with separate read and write
cursors.

The channel interface provides two operations:
\begin{itemize}[leftmargin=1.4em]
  \item \syscall{chan\_send(chan\_cap, msg)}: enqueue a message. Returns
    immediately; fails if the ring is full.
  \item \syscall{chan\_recv(chan\_cap)}: dequeue a message. Returns
    immediately; fails if the ring is empty.
\end{itemize}

\section{Service Registry}

The kernel maintains a name-to-endpoint map that allows services to discover
each other:

\begin{itemize}[leftmargin=1.4em]
  \item \syscall{SvcRegister} (syscall~130): register a string name with
    an endpoint capability.
  \item \syscall{SvcLookup} (syscall~131): look up a name and receive a
    derived endpoint capability.
\end{itemize}

\section{IPC Data Layout}

IPC messages carry data in registers with the following layout constraints:

\begin{itemize}[leftmargin=1.4em]
  \item \textbf{Send}: up to 40 bytes of inline data (registers 3--7,
    5 registers at 8 bytes each).
  \item \textbf{Receive}: up to 64 bytes (registers 0--7, 8 registers at
    8 bytes each), with the byte count in the tag field.
  \item For larger transfers, the \LUCAS{} layer uses a multi-chunk IPC
    loop that reassembles data 64 bytes at a time (used for TCP/HTTP).
\end{itemize}


% ------------------------------------------------------------------------------
\chapter{Scheduler}
\label{ch:sched}
% ------------------------------------------------------------------------------

\section{Multi-Level Queue Design}

The scheduler implements a four-level priority MLQ with the following queues:

\begin{table}[H]
\centering
\caption{Scheduler priority levels.}
\begin{tabular}{c l l l}
\toprule
\textbf{Level} & \textbf{Name} & \textbf{Use Case} & \textbf{Policy} \\
\midrule
0 & Realtime & IRQ handlers, timer ISR & Fixed priority, no preemption \\
1 & High     & VMM, FS server          & Round-robin, short quantum \\
2 & Normal   & Applications            & Round-robin, standard quantum \\
3 & Low      & Background tasks        & Best-effort, long quantum \\
\bottomrule
\end{tabular}
\end{table}

\section{Work Stealing}

Each CPU maintains a Chase-Lev work-stealing deque
(\file{kernel/src/sched/wsdeque.rs}). When a CPU's local run queue is empty,
it attempts to steal threads from other CPUs' deques. This provides
automatic load balancing across cores without centralized coordination.

The Chase-Lev deque provides:
\begin{itemize}[leftmargin=1.4em]
  \item O(1) local push/pop (LIFO for cache locality).
  \item O(1) remote steal (FIFO for fairness).
  \item Lock-free operation using atomic compare-and-swap.
\end{itemize}

\section{Resource Limits}

Each process has configurable resource limits enforced by the scheduler
and frame allocator:
\begin{itemize}[leftmargin=1.4em]
  \item \textbf{CPU ticks}: maximum number of scheduler ticks before
    the process is descheduled. Enforced via \syscall{SYS\_RESOURCE\_LIMIT}
    (syscall~141).
  \item \textbf{Memory pages}: maximum number of physical frames a process
    may allocate. Enforced in the frame allocator.
\end{itemize}

\section{SMP Support}

The kernel supports up to 16 CPUs. Each Application Processor (AP) is woken
via the standard INIT-SIPI-SIPI IPI sequence. Key SMP details:

\begin{itemize}[leftmargin=1.4em]
  \item Each CPU has its own LAPIC timer for preemption.
  \item Cross-core wakeup uses IPIs (Inter-Processor Interrupts).
  \item Per-CPU data is accessed via a per-CPU segment or the LAPIC ID.
  \item SSE must be enabled on \emph{all} CPUs (BSP and APs), as musl and
    Rust binaries use SSE instructions (e.g., \texttt{movaps}).
  \item The LAPIC timer is calibrated using the PIT (Programmable Interval
    Timer) with a CPUID.15H/TSC fallback chain for hypervisors.
\end{itemize}


% ------------------------------------------------------------------------------
\chapter{Security}
\label{ch:security}
% ------------------------------------------------------------------------------

\section{Threat Model}

\sotBSD{} defends against:
\begin{itemize}[leftmargin=1.4em]
  \item Compromised userspace services (no ambient authority).
  \item Capability confusion attacks (generation counters).
  \item Stack-based exploits (guard pages, canaries, W\^{}X).
  \item Address space layout attacks (ASLR).
  \item Adversarial reconnaissance (deception engine, Chapter~\ref{ch:deception}).
\end{itemize}

\section{W\^{}X Enforcement}

The kernel enforces Write XOR Execute at the page table level. No page may
have both the writable (W) and executable (X) bits set simultaneously. The
\syscall{SYS\_PROTECT} syscall rejects any request that violates this
invariant.

\section{Stack Security}

Three mechanisms protect stacks:
\begin{description}[leftmargin=1.4em]
  \item[Guard pages.] An unmapped page below each stack triggers \#PF on
    overflow.
  \item[Stack canaries.] A fixed compile-time sentinel value is placed
    between the return address and local variables. The kernel checks this
    value on context switch.
  \item[ASLR.] Stack base addresses are randomized within a configurable
    range.
\end{description}

\section{Signed Boot Chain}

Every binary in the initrd CPIO archive is Ed25519-signed. The public key
is compiled into the init service via
\file{services/init/src/sigkey\_generated.rs}. At boot, each binary's SHA-256
hash is verified against its Ed25519 signature before loading. Public-key
pinning rejects swap-with-self-signed attempts.

\section{Provenance Tracking}

The SOT provenance subsystem maintains a SHA-256 hash-linked DAG of all
object mutations. Each CPU has a dedicated 4096-entry SPSC ring buffer that
records provenance events. The hash chain ensures that any tampering with
the provenance log is detectable.


% ==============================================================================
% PART III --- PERSONALITIES
% ==============================================================================
\part{Personalities}

% ------------------------------------------------------------------------------
\chapter{LUCAS: Linux ABI Compatibility}
\label{ch:lucas}
% ------------------------------------------------------------------------------

\section{Architecture}

\LUCAS{} (Linux User-mode Compatibility And Syscall translation) implements
a comprehensive Linux system call interface on top of the \sotBSD{} kernel.
Rather than emulating a Linux kernel, \LUCAS{} translates Linux syscalls
into sequences of \sotBSD{} IPC calls and local operations within the init
service's \texttt{child\_handler}.

The key architectural insight is that the kernel intercepts
\texttt{arch\_prctl} at the syscall level (for TLS/FS\_BASE management) and
redirects all other Linux syscalls via IPC to the init service, which
handles them in the context of the calling process.

\section{Supported Syscalls}

\LUCAS{} implements 127+ Linux syscall match arms in the \texttt{child\_handler}
function. Table~\ref{tab:lucas-syscalls} lists the major categories.

\begin{longtable}{l p{9cm}}
\caption{LUCAS Linux syscall categories.} \label{tab:lucas-syscalls} \\
\toprule
\textbf{Category} & \textbf{Syscalls} \\
\midrule
\endfirsthead
\toprule
\textbf{Category} & \textbf{Syscalls} \\
\midrule
\endhead
\midrule
\multicolumn{2}{r}{\textit{Continued on next page}} \\
\endfoot
\bottomrule
\endlastfoot
File I/O & open, openat, read, write, close, lseek, pread64, pwrite64,
           readv, writev, preadv, pwritev, dup, dup2, dup3, fcntl \\
Directory & mkdir, mkdirat, rmdir, unlink, unlinkat, rename, renameat,
            renameat2, getdents64 \\
File metadata & stat, fstat, fstatat, lstat, statfs, fstatfs, access,
                faccessat, faccessat2, readlink, readlinkat \\
Process & fork (clone CoW), execve, exit, exit\_group, wait4, waitid,
          getpid, gettid, getppid, clone (CLONE\_THREAD) \\
Memory & brk, mmap, munmap, mprotect, mremap, madvise \\
Signals & rt\_sigaction, rt\_sigprocmask, rt\_sigreturn, rt\_sigpending,
          rt\_sigsuspend, rt\_sigtimedwait, tgkill, sigaltstack \\
Sockets & socket, bind, connect, send, recv, sendto, recvfrom, recvmsg,
          poll, epoll\_create1, epoll\_ctl, epoll\_wait, socketpair \\
Threading & clone (CLONE\_VM/FILES/SIGHAND/THREAD), futex, set\_tid\_address,
            gettid \\
Time & clock\_gettime, clock\_getres, clock\_nanosleep, gettimeofday, time,
       nanosleep \\
Identity & getuid, getgid, geteuid, getegid, getresuid, getresgid, setuid,
           setgid, getgroups, setgroups \\
IPC & pipe, pipe2, eventfd, timerfd\_create, timerfd\_settime, signalfd,
      shmget, shmat, shmctl \\
System & uname, sysinfo, prctl, arch\_prctl, personality, getrandom,
         prlimit64, getrusage \\
File ops & ftruncate, fsync, fdatasync, fallocate, flock \\
Special & ioctl, select, pselect6, memfd\_create, inotify\_init1 \\
\end{longtable}

\section{Process Model}

\LUCAS{} supports the full Unix process lifecycle:

\subsection{Fork}

\texttt{fork()} is implemented via \syscall{SYS\_AS\_CLONE} (syscall~125),
which creates a real copy-on-write address space clone. Three critical bugs
were fixed in the fork implementation:

\begin{enumerate}[leftmargin=1.4em]
  \item \textbf{Stack save/restore}: The parent's 4\,KiB stack at
    \texttt{fork\_rsp} is saved to a static buffer (\texttt{FORK\_STACK\_BUF})
    before entering the fork-child loop, and restored after exec/exit. This
    prevents the vfork child from corrupting the parent's stack.
  \item \textbf{FD state restoration}: Parent FD state is always restored
    after the fork-child phase completes.
  \item \textbf{Pipe reference counting}: \texttt{PIPE\_WRITE\_REFS} and
    \texttt{PIPE\_READ\_REFS} arrays prevent premature EOF from child close.
\end{enumerate}

\subsection{Exec}

\texttt{execve()} loads ELF binaries from either the initrd CPIO archive or
the VFS. The \texttt{exec\_loaded\_elf()} function in \file{services/init/src/exec.rs}
is the shared ELF loader that handles both static and dynamically-linked
binaries:

\begin{itemize}[leftmargin=1.4em]
  \item Static musl binaries: loaded directly at their ELF load address.
  \item Dynamic binaries: the PT\_INTERP segment is read to find the
    interpreter (e.g., \texttt{ld-musl-x86\_64.so.1}), which is loaded at
    \texttt{INTERP\_LOAD\_BASE} (\hex{6000000}). The auxiliary vector
    includes \texttt{AT\_BASE} so the linker knows its own load address.
\end{itemize}

\subsection{Thread Groups}

\LUCAS{} implements thread groups for multi-threaded applications:

\begin{itemize}[leftmargin=1.4em]
  \item \texttt{PROC\_TGID}: maps threads to their thread group leader.
  \item \texttt{PROC\_FD\_GROUP}: threads in the same group share file
    descriptors.
  \item \texttt{PROC\_MEM\_GROUP}: threads share the address space.
  \item \texttt{PROC\_SIG\_GROUP}: threads share signal handlers.
\end{itemize}

\texttt{clone()} with \texttt{CLONE\_THREAD | CLONE\_VM | CLONE\_FILES |
CLONE\_SIGHAND} creates a new thread in the same group, using a clone
trampoline that sets up TLS via \texttt{arch\_prctl}.

\section{vDSO}

The virtual Dynamic Shared Object (vDSO) is a minimal ELF shared library
(\texttt{ET\_DYN}) injected as a single 4\,KiB page at \hex{B80000}. It
exports:

\begin{itemize}[leftmargin=1.4em]
  \item \texttt{\_\_vdso\_clock\_gettime}: RDTSC-based time, assuming 2\,GHz.
  \item Stubs for \texttt{gettimeofday}, \texttt{time}, \texttt{getcpu}.
  \item Fork-restore trampoline at \hex{B80140}.
  \item Signal trampoline at \hex{B80700}.
  \item Signal return stub at \hex{B80710}.
  \item Exit stub at \hex{B80730}.
\end{itemize}

The vDSO includes a full ELF structure: \texttt{PT\_LOAD}, \texttt{PT\_DYNAMIC},
\texttt{.dynsym}, \texttt{.dynstr}, SysV \texttt{.hash},
\texttt{.gnu.version}, and \texttt{.gnu.version\_d} with the
\texttt{LINUX\_2.6} version tag. \texttt{AT\_SYSINFO\_EHDR} (auxv type~33)
points every child process to the vDSO.

\section{/proc Emulation}

\texttt{child\_handler} intercepts \texttt{openat} calls to \texttt{/proc/}
and \texttt{/sys/} paths, returning virtual file descriptors (kind=15) whose
contents are generated on-the-fly. Supported paths include:

\begin{multicols}{2}
\begin{itemize}[leftmargin=1em,itemsep=0pt]
  \item \texttt{/proc/self/maps}
  \item \texttt{/proc/self/status}
  \item \texttt{/proc/self/stat}
  \item \texttt{/proc/self/auxv}
  \item \texttt{/proc/self/environ}
  \item \texttt{/proc/self/exe} (readlinkat)
  \item \texttt{/proc/cpuinfo}
  \item \texttt{/proc/meminfo}
  \item \texttt{/proc/N/status}
  \item \texttt{/sys/devices/system/cpu/online}
\end{itemize}
\end{multicols}


% ------------------------------------------------------------------------------
\chapter{BSD and rump Kernel Personality}
\label{ch:bsd}
% ------------------------------------------------------------------------------

\section{rump Kernel Integration}

The BSD personality uses a real NetBSD rump kernel linked into a freestanding
\sotOS{} service. The rump kernel provides a complete NetBSD kernel
environment including VFS, device drivers, and protocol stacks, running
entirely in userspace.

The integration comprises:
\begin{itemize}[leftmargin=1.4em]
  \item \textbf{Libraries:} \texttt{librump.a}, \texttt{librumpvfs.a},
    \texttt{librumpfs\_ffs.a}, \texttt{librumpdev.a},
    \texttt{librumpdev\_disk.a} (approximately 2\,MiB total).
  \item \textbf{Hypercalls:} 60 hand-rolled \texttt{rumpuser\_*} functions
    in \texttt{rumpuser\_sot.c} that bridge the rump kernel to \sotBSD{}
    IPC and memory management.
  \item \textbf{Block device:} Connected to the virtio-blk or NVMe driver
    via IPC for persistent storage.
\end{itemize}

\section{Illumos Heirlooms}

The v0.2.0+ release includes three Illumos-inspired subsystems:

\subsection{SMF: Service Management Facility}

SMF provides dependency-tracked service management with automatic respawn:

\begin{itemize}[leftmargin=1.4em]
  \item Services declare dependencies as a topological DAG.
  \item The SMF manager starts services in dependency order.
  \item Failed services are automatically restarted via
    \syscall{SYS\_THREAD\_NOTIFY} (syscall~143).
  \item Cascade failures propagate through the dependency graph.
\end{itemize}

\subsection{Crossbow: Virtual Networking}

Project Crossbow provides:
\begin{itemize}[leftmargin=1.4em]
  \item O(1) VNIC provisioning and revocation.
  \item Per-port token-bucket rate limiting.
  \item Locally-administered MAC address generation.
  \item Virtual switch with multiple forwarding modes.
\end{itemize}

\subsection{FMA: Fault Management Architecture}

The Fault Management Architecture provides predictive self-healing across
five scenarios:

\begin{enumerate}[leftmargin=1.4em]
  \item Disk retry threshold exceeded.
  \item ECC correctable memory errors above threshold.
  \item ECC uncorrectable memory errors (immediate response).
  \item Thermal events (CPU throttling or shutdown).
  \item PCI link degradation (lane width reduction).
\end{enumerate}

Six prioritized correlation rules process events from disk, memory, thermal,
and PCI subsystems, combined with provenance anomaly counters.


% ------------------------------------------------------------------------------
\chapter{Deception Engine}
\label{ch:deception}
% ------------------------------------------------------------------------------

\section{Capability Interposition Pipeline}

The deception engine operates through capability interposition: when a
syscall traverses a capability marked with a deception policy, the engine
transforms the request or response according to the active deception
profile.

The pipeline processes four stages:
\begin{enumerate}[leftmargin=1.4em]
  \item \textbf{Intercept}: The capability's interposition policy
    (Section~\ref{ch:caps}) triggers the deception handler.
  \item \textbf{Profile lookup}: The active deception profile determines
    what fabrication to apply.
  \item \textbf{Transform}: The response is rewritten according to the
    profile's rules.
  \item \textbf{Consistency check}: The consistency engine verifies
    cross-channel coherence before delivering the result.
\end{enumerate}

\section{Deception Profiles}

The engine includes five built-in profiles:

\begin{table}[H]
\centering
\caption{Built-in deception profiles.}
\begin{tabularx}{\textwidth}{l X}
\toprule
\textbf{Profile} & \textbf{Description} \\
\midrule
Ubuntu 22.04  & Emulates a standard Ubuntu webserver with Apache \\
Debian        & Mimics Debian stable with default packages \\
CentOS        & Presents as a CentOS/RHEL database server \\
FreeBSD       & Emulates a FreeBSD server \\
Alpine        & Lightweight Alpine Linux container \\
\bottomrule
\end{tabularx}
\end{table}

Each profile defines consistent responses for:
\begin{itemize}[leftmargin=1.4em]
  \item \texttt{uname} syscall (kernel version, hostname, architecture).
  \item \texttt{/proc/version} and \texttt{/proc/cpuinfo} contents.
  \item \texttt{/etc/os-release} and \texttt{/etc/hostname}.
  \item File system layout and installed packages.
  \item Network service banners (SSH, HTTP, etc.).
\end{itemize}

\section{Live Migration}

The deception engine supports live profile migration with the following
guarantees:

\begin{itemize}[leftmargin=1.4em]
  \item \textbf{Atomic CSpace swap}: The entire capability space is swapped
    atomically, ensuring no inconsistent intermediate state.
  \item \textbf{Sub-10\,$\mu$s latency}: The swap is performed in under
    10 microseconds, making it undetectable by timing analysis.
  \item \textbf{Connection preservation}: Active network connections are
    maintained across profile switches.
\end{itemize}

\section{Consistency Engine}

The consistency engine ensures cross-channel coherence. When the deception
profile claims to be ``Ubuntu 22.04'', all channels must agree:

\begin{itemize}[leftmargin=1.4em]
  \item \texttt{uname -r} returns the correct kernel version.
  \item \texttt{/etc/os-release} matches.
  \item \texttt{/proc/version} is consistent.
  \item Network banners use appropriate version strings.
  \item File timestamps and package lists are plausible.
\end{itemize}

\section{Detection Resistance}

To resist timing-based detection, the engine employs:
\begin{itemize}[leftmargin=1.4em]
  \item Calibrated delays matching the expected profile's syscall latency.
  \item Random jitter added to response times.
  \item Cache-line isolation to prevent microarchitectural side channels.
\end{itemize}


% ==============================================================================
% PART IV --- SERVICES AND DRIVERS
% ==============================================================================
\part{Services \& Drivers}

% ------------------------------------------------------------------------------
\chapter{sotFS: The Filesystem}
\label{ch:sotfs}
% ------------------------------------------------------------------------------

\section{Overview}

\sotFS{} is the native filesystem for \sotBSD{}, designed around a type graph
model with formal verification guarantees. The implementation consists of 7
Rust crates, 76 tests, and a mountable FUSE prototype for Linux development.

\section{Type Graph Model}

The filesystem state is represented as a directed graph $TG = (V, E)$ where
nodes and edges are typed. The type graph has 6 node types and 6 edge types.

\begin{figure}[H]
\centering
\begin{tikzpicture}[
  nodetype/.style={draw=tnblue, thick, fill=tnblue!10, circle,
                   minimum size=1.4cm, font=\sffamily\small, align=center},
  edge/.style={-{Stealth[length=2.5mm]}, thick},
]
  % Nodes
  \node[nodetype] (inode) at (0,3)    {Inode};
  \node[nodetype] (dir)   at (4,3)    {Dir};
  \node[nodetype, fill=tngreen!15] (cap) at (8,3) {Cap};
  \node[nodetype, fill=tnyellow!15] (tx) at (0,0) {Tx};
  \node[nodetype, fill=tnpurple!15] (ver) at (4,0) {Version};
  \node[nodetype, fill=tnorange!15] (blk) at (8,0) {Block};

  % Edges
  \draw[edge, tnblue] (dir) -- (inode)
    node[midway, above, font=\scriptsize] {contains};
  \draw[edge, tngreen] (cap) -- (inode)
    node[midway, above, font=\scriptsize] {authorizes};
  \draw[edge, tnred] (dir) to[bend left=20] (dir)
    node[midway, right, font=\scriptsize] {parent};
  \draw[edge, tnyellow] (tx) -- (inode)
    node[midway, left, font=\scriptsize] {modifies};
  \draw[edge, tnpurple] (inode) -- (ver)
    node[midway, above, font=\scriptsize] {has\_version};
  \draw[edge, tnorange] (ver) -- (blk)
    node[midway, above, font=\scriptsize] {maps\_to};

  % Legend
  \node[font=\sffamily\scriptsize\color{tncomment}, anchor=north west]
    at (-1.5,-1.5) {
    \begin{tabular}{@{}ll@{}}
    \textbf{Node types:} & Inode, Directory, Capability, Transaction, Version, Block \\
    \textbf{Edge types:} & contains, authorizes, parent, modifies, has\_version, maps\_to \\
    \end{tabular}};
\end{tikzpicture}
\caption{sotFS type graph with 6 node types and 6 edge types. Edges show
         the relationships between filesystem entities.}
\label{fig:sotfs-typegraph}
\end{figure}

\section{DPO Rewriting Rules}

POSIX filesystem operations are modeled as Double Pushout (DPO) graph
rewriting rules. Each rule specifies a left-hand side (pattern to match),
a right-hand side (replacement), and an interface (preserved context).

\begin{definition}[DPO Rule for CREATE]
\label{def:dpo-create}
Given a directory node $d$ and a fresh inode $i$, the CREATE operation is
the DPO rule:
\[
  L \xleftarrow{l} K \xrightarrow{r} R
\]
where:
\begin{itemize}[leftmargin=1.4em]
  \item $L = \{d : \text{Dir}\}$ --- the left-hand side contains only the
    parent directory.
  \item $K = \{d : \text{Dir}\}$ --- the interface preserves the directory.
  \item $R = \{d : \text{Dir},\; i : \text{Inode},\;
    (d \xrightarrow{\text{contains}} i)\}$ --- the right-hand side adds
    the new inode and a ``contains'' edge.
\end{itemize}
The rule is applied in context $G$ (the current filesystem state) to produce
$H$ (the new state): $G \Rightarrow_{\text{CREATE}} H$.
\end{definition}

Table~\ref{tab:dpo-rules} lists all DPO rules.

\begin{table}[H]
\centering
\caption{POSIX operations as DPO rewriting rules.}
\label{tab:dpo-rules}
\begin{tabularx}{\textwidth}{l X}
\toprule
\textbf{VfsOp} & \textbf{DPO Rule Description} \\
\midrule
CREATE  & Add Inode node, add ``contains'' edge from Dir \\
MKDIR   & Add Dir node, add ``contains'' and ``parent'' edges \\
RMDIR   & Remove Dir node (requires empty), remove edges \\
LINK    & Add ``contains'' edge from target Dir to existing Inode \\
UNLINK  & Remove ``contains'' edge (remove Inode if link count = 0) \\
RENAME  & Remove ``contains'' from source Dir, add to target Dir \\
WRITE   & Add Version node, add ``has\_version'' edge, add Block nodes \\
\bottomrule
\end{tabularx}
\end{table}

\section{Bounded Treewidth}

\begin{theorem}[Bounded Treewidth]
The treewidth of the sotFS type graph $TG$ satisfies:
\[
  \mathrm{tw}(TG) \leq \mathrm{LINK\_MAX} + O(1)
\]
Without hard links, the bound tightens to $\mathrm{tw}(TG) \leq 6$.
\end{theorem}

This bounded treewidth guarantee ensures that graph algorithms on the
filesystem state (such as anomaly detection) run in linear time.

\section{Ollivier-Ricci Curvature}

\sotFS{} uses Ollivier-Ricci curvature as an anomaly detection metric. The
discrete curvature of an edge $(u, v)$ measures how well-connected the
neighborhoods of $u$ and $v$ are. Anomalous filesystem activity (e.g.,
ransomware creating many files in a flat directory) produces edges with
strongly negative curvature, triggering alerts.

\section{Deceptive Subgraph Projections}

The deception functor $D_{FS} : \mathbf{TG} \to \mathbf{TG}$ produces a
projected view of the filesystem for deception targets. This functor can:
\begin{itemize}[leftmargin=1.4em]
  \item Hide real files while presenting fabricated directory listings.
  \item Inject honeypot files that trigger alerts when accessed.
  \item Present a plausible filesystem layout matching the active deception
    profile.
\end{itemize}

\section{Implementation Quality}

The \sotFS{} implementation has been verified through:
\begin{itemize}[leftmargin=1.4em]
  \item \textbf{4 TLA+ specifications} checked with TLC.
  \item \textbf{Typestate enforcement} at compile time.
  \item \textbf{loom concurrency tests} for lock-free data structures.
  \item \textbf{proptest}: property-based testing that found 2 bugs.
  \item \textbf{cargo-fuzz}: 300K runs with 0 crashes.
  \item \textbf{criterion}: performance benchmarks for all operations.
  \item \textbf{76 unit and integration tests}.
  \item \textbf{FUSE prototype}: mountable on Linux for real-world testing.
\end{itemize}


% ------------------------------------------------------------------------------
\chapter{Device Drivers}
\label{ch:drivers}
% ------------------------------------------------------------------------------

\section{Driver Architecture}

All device drivers run in Ring~3 as capability-holding services. The kernel
provides IRQ virtualization: hardware interrupts are demultiplexed and
delivered as IPC messages to registered handlers. This architecture ensures
that a buggy or compromised driver cannot crash the kernel.

\section{PCI Enumeration}

The \texttt{sotos-pci} library provides PCI bus enumeration. It scans the
configuration space to discover devices and their BARs (Base Address
Registers), then maps MMIO regions into the driver's address space with
UC (Uncacheable) flags.

\section{virtio Drivers}

\subsection{virtio-blk}

The virtio-blk driver provides block device access through virtio's split
virtqueue mechanism. It handles read and write requests from the filesystem
and other services via IPC.

\subsection{virtio-net}

The virtio-net driver provides network interface access. It integrates with
the networking stack (Chapter~\ref{ch:net}) to send and receive Ethernet
frames.

\section{NVMe Driver}

The NVMe driver (\texttt{sotos-nvme}) communicates with NVMe SSDs using:
\begin{itemize}[leftmargin=1.4em]
  \item Admin and I/O submission/completion queue pairs.
  \item PRP (Physical Region Page) lists for scatter-gather DMA.
  \item IPC-based interface for filesystem requests.
\end{itemize}

\section{xHCI USB Driver}

The xHCI driver (\texttt{sotos-xhci}) manages USB 3.x host controllers:
\begin{itemize}[leftmargin=1.4em]
  \item Device slot allocation and endpoint configuration.
  \item Transfer ring management (command, event, transfer rings).
  \item HID keyboard support: USB keyboard events are translated and
    delivered to the keyboard ring buffer.
\end{itemize}

\section{PS/2 Keyboard and Mouse}

The PS/2 driver handles legacy keyboard and mouse input:
\begin{itemize}[leftmargin=1.4em]
  \item Keyboard: scan code set 1/2 translation, key-up/key-down events
    written to the ring buffer at \hex{510000}.
  \item Mouse: PS/2 mouse protocol with movement and button events
    written to the ring buffer at \hex{520000}.
\end{itemize}

\section{AC97 Audio}

The AC97 audio driver provides basic audio output through the AC97 codec
interface, supporting PCM playback at standard sample rates.

\section{Minimal Service Skeleton}

Listing~\ref{lst:service-skeleton} shows the minimal code to create a
new service that registers with the service registry and enters an IPC
loop.

\begin{lstlisting}[caption={Minimal service skeleton.},label={lst:service-skeleton}]
#![no_std]
#![no_main]

use sotos_common::sys;
use sotos_common::ipc::Message;

#[no_mangle]
pub extern "C" fn _start() -> ! {
    // Read our capabilities from BootInfo
    let boot_info = unsafe { &*(0xB00000 as *const BootInfo) };
    let root_cap = boot_info.root_cap;

    // Create an endpoint for receiving requests
    let ep = sys::endpoint_create()
        .expect("endpoint_create failed");

    // Register our service name so others can find us
    sys::svc_register("my_service", ep)
        .expect("svc_register failed");

    // Main IPC loop: receive requests and send replies
    loop {
        let mut request = Message::new();
        let badge = sys::recv(ep, &mut request);

        // Dispatch based on the message tag
        let mut reply = Message::new();
        match request.tag() {
            1 => { /* handle command 1 */ }
            2 => { /* handle command 2 */ }
            _ => { reply.set_error(-1); }
        }

        sys::reply(badge, &reply);
    }
}
\end{lstlisting}


% ------------------------------------------------------------------------------
\chapter{Networking Stack}
\label{ch:net}
% ------------------------------------------------------------------------------

\section{Protocol Stack}

The networking stack is implemented as a userspace service with the following
layers:

\begin{description}[leftmargin=1.4em]
  \item[Layer 2: Ethernet.] Frame transmission and reception via
    virtio-net driver.
  \item[ARP.] Address Resolution Protocol with dynamic cache.
  \item[DHCP.] Dynamic host configuration with random transaction IDs.
  \item[Layer 3: IP.] IPv4 packet routing and fragmentation.
  \item[ICMP.] Ping echo request/reply and error messages.
  \item[Layer 4: UDP.] Connectionless datagram transport. Used for DNS.
  \item[Layer 4: TCP.] Connection-oriented transport with:
    \begin{itemize}[leftmargin=1em]
      \item Three-way handshake (SYN, SYN-ACK, ACK).
      \item Retransmission with exponential backoff.
      \item RST handling for connection reset.
      \item TCP Reno congestion control.
      \item 16\,KiB receive buffer.
    \end{itemize}
  \item[DNS.] Domain name resolution via UDP to configured resolver
    (default: \texttt{10.0.2.3} under QEMU SLIRP).
\end{description}

\section{Socket API}

\LUCAS{} exposes a BSD socket API to Linux binaries:

\begin{itemize}[leftmargin=1.4em]
  \item \texttt{socket(AF\_INET, SOCK\_STREAM|SOCK\_DGRAM, 0)}: creates
    TCP or UDP socket FDs.
  \item \texttt{bind}, \texttt{connect}, \texttt{listen}, \texttt{accept}.
  \item \texttt{send}, \texttt{recv}, \texttt{sendto}, \texttt{recvfrom},
    \texttt{recvmsg}.
  \item \texttt{poll}, \texttt{epoll\_create1}, \texttt{epoll\_ctl},
    \texttt{epoll\_wait}.
\end{itemize}

The socket proxy translates these into IPC messages to the net service:
\begin{itemize}[leftmargin=1.4em]
  \item TCP send: \texttt{regs=[conn\_id, 0, data\_len, data...]} (data
    in regs[3..], 40 bytes max per message).
  \item UDP sendto: \texttt{regs=[dst\_ip, dst\_port, src\_port, data\_len,
    data...]} (data in regs[4..]).
  \item Receive: reply tag = byte count, data in regs[0..7] (64 bytes max).
\end{itemize}

\section{HTTP Client}

The networking stack supports HTTP/1.1 GET requests sufficient for tools
like \texttt{busybox wget}. Multi-chunk IPC assembly handles responses
larger than 64 bytes by looping the receive call.


% ==============================================================================
% PART V --- OPERATIONS
% ==============================================================================
\part{Operations}

% ------------------------------------------------------------------------------
\chapter{Building and Running}
\label{ch:build}
% ------------------------------------------------------------------------------

\section{Prerequisites}

\begin{itemize}[leftmargin=1.4em]
  \item Rust nightly toolchain (for \texttt{build-std} and
    \texttt{abi\_x86\_interrupt}).
  \item \texttt{x86\_64-unknown-none} target:
    \texttt{rustup target add x86\_64-unknown-none}.
  \item Python 3 (invoked as \texttt{python} on Windows).
  \item QEMU (at \texttt{C:\textbackslash Program Files\textbackslash qemu\textbackslash}
    on Windows).
  \item \texttt{just} command runner.
\end{itemize}

\section{Build Commands}

\begin{table}[H]
\centering
\caption{Build and run commands.}
\begin{tabularx}{\textwidth}{l X}
\toprule
\textbf{Command} & \textbf{Description} \\
\midrule
\texttt{just build}     & Compile kernel and all services \\
\texttt{just run}       & Build + create disk image + boot QEMU (serial) \\
\texttt{just run-gui}   & Same as \texttt{run} but with QEMU display window \\
\texttt{just run-net}   & Boot with virtio-net enabled \\
\texttt{just run-nvme}  & Boot with NVMe device \\
\texttt{just run-xhci}  & Boot with xHCI USB controller \\
\texttt{just run-fast}  & Boot with WHPX acceleration (${\sim}$25s to Tier~5) \\
\texttt{just debug}     & Boot with GDB server (\texttt{target remote :1234}) \\
\bottomrule
\end{tabularx}
\end{table}

\section{Manual Build Steps}

\begin{lstlisting}[language={},caption={Manual build and boot sequence.}]
# Compile the kernel
cargo build --package sotos-kernel

# Create the disk image (GPT + FAT32 with LFN)
python scripts/mkimage.py

# Boot in QEMU (2048M RAM, serial to terminal)
"C:/Program Files/qemu/qemu-system-x86_64.exe" \
  -drive format=raw,file=target/sotos.img \
  -serial stdio -display none -no-reboot -m 2048M
\end{lstlisting}

\section{Disk Image}

The disk image is built by \file{scripts/mkimage.py}, a pure-Python GPT +
FAT32 image builder with Long File Name (LFN) support. It creates a raw
disk image containing:

\begin{itemize}[leftmargin=1.4em]
  \item A GPT partition table with an EFI System Partition.
  \item The Limine bootloader (BIOS stage 1 and 2).
  \item The kernel ELF binary.
  \item The initrd CPIO archive containing all services.
  \item The \texttt{limine.conf} boot configuration.
\end{itemize}


% ------------------------------------------------------------------------------
\chapter{Debugging}
\label{ch:debug}
% ------------------------------------------------------------------------------

\section{Serial Output}

All kernel debug output goes through COM1 serial (\texttt{kprintln!()} macro).
QEMU's \texttt{-serial stdio} flag redirects this to the terminal. The
\texttt{kdebug!()} macro is gated behind the \texttt{verbose} feature flag,
allowing conditional debug output.

\textbf{Critical invariant}: Serial must be initialized \emph{first} in
\texttt{kmain()}, before any assertions or debug prints. A panic before
serial initialization produces no diagnostics.

\section{GDB Debugging}

\begin{lstlisting}[language={},caption={GDB debugging session.}]
# Terminal 1: boot with GDB server
just debug

# Terminal 2: connect GDB
gdb -ex "target remote :1234"
(gdb) break kmain
(gdb) continue
\end{lstlisting}

\section{Common Gotchas}

\begin{description}[leftmargin=1.4em]
  \item[ET\_DYN rejection.] Limine rejects dynamically-linked ELF files.
    Ensure \texttt{-C relocation-model=static} is set.
  \item[Cargo config merging.] Child crate builds inherit parent
    \texttt{.cargo/config.toml}. Use \texttt{CARGO\_ENCODED\_RUSTFLAGS}
    to override.
  \item[TCG spin deadlock.] Under QEMU TCG (software emulation), tight
    spin loops may deadlock. Use \texttt{yield} or \texttt{hlt}.
  \item[WHPX strict mode.] WHPX (Windows Hypervisor Platform) enforces
    strict x86 semantics (FXRSTOR reserved bits, MXCSR validation) that
    TCG ignores.
  \item[IST and \#PF.] Never assign an IST entry to \#PF---shared IST
    stacks corrupt under concurrent CoW faults.
  \item[Init stack size.] Use 4 pages (16\,KiB), not 1. Smaller stacks
    overflow during service initialization.
  \item[MMIO UC flags.] Device MMIO pages must be mapped with Uncacheable
    flags to prevent CPU write-back caching.
  \item[compiler\_builtins memcmp.] In debug builds, \texttt{compiler\_builtins}
    \texttt{memcmp} triggers UB checks. Use byte-by-byte comparison in
    \file{kernel/src/initrd.rs}.
  \item[Atomic thread redirect.] Use \texttt{spawn\_user\_with\_redirect()}
    ---setting redirect\_ep after thread creation races with early syscalls.
\end{description}


% ------------------------------------------------------------------------------
\chapter{Syscall Reference}
\label{ch:syscalls}
% ------------------------------------------------------------------------------

\section{Kernel Syscalls}

Table~\ref{tab:kern-syscalls} lists all kernel-native syscalls.

\begin{longtable}{c l p{3.5cm} p{5cm}}
\caption{Kernel syscall reference.} \label{tab:kern-syscalls} \\
\toprule
\textbf{\#} & \textbf{Name} & \textbf{Arguments} & \textbf{Returns / Description} \\
\midrule
\endfirsthead
\toprule
\textbf{\#} & \textbf{Name} & \textbf{Arguments} & \textbf{Returns / Description} \\
\midrule
\endhead
\midrule
\multicolumn{4}{r}{\textit{Continued on next page}} \\
\endfoot
\bottomrule
\endlastfoot
24  & UnmapFree      & vaddr          & Unmap page + free frame \\
130 & SvcRegister    & name, ep\_cap   & Register service name \\
131 & SvcLookup      & name           & Returns derived ep cap \\
132 & InitrdRead     & filename, buf  & Read file from initrd CPIO \\
133 & BootInfoWrite  & as\_cap, addr   & Write BootInfo into target AS \\
134 & Protect        & addr, len, prot & mprotect with W\^{}X \\
135 & CallTimeout    & cap, msg, ticks & IPC call with timeout \\
140 & ThreadInfo     &                & Get tid/priority/ticks/mem \\
141 & ResourceLimit  & ticks, pages   & Set CPU/memory limits \\
142 & ThreadCount    &                & Get total thread count \\
143 & ThreadNotify   & tid            & Wake thread (for SMF) \\
160 & SetFsBase      & addr           & Set FS\_BASE MSR (TLS) \\
161 & GetFsBase      &                & Get FS\_BASE MSR \\
172 & GetThreadRegs  & raw\_tid        & Read saved registers \\
173 & SignalEntry    & thread\_cap, addr & Set signal trampoline \\
174 & SignalInject   & raw\_tid, signum & Set pending signal bit \\
252 & DebugFreeFrames &               & Return free frame count \\
\end{longtable}

\section{SOT Syscalls}

See Table~\ref{tab:sot-syscalls} in Chapter~\ref{ch:arch} for the 12 SOT
exokernel syscalls (numbers 300--311).

\section{FD Kind Codes}

Table~\ref{tab:fd-kinds} lists the file descriptor kind codes used by
\LUCAS{} to dispatch operations to the correct handler.

\begin{table}[H]
\centering
\caption{File descriptor kind codes.}
\label{tab:fd-kinds}
\begin{tabular}{c l l}
\toprule
\textbf{Kind} & \textbf{Name} & \textbf{Description} \\
\midrule
1  & Stdin      & Standard input (keyboard ring) \\
2  & Stdout     & Standard output (serial) \\
3  & Stderr     & Standard error (serial) \\
8  & Null       & \texttt{/dev/null} \\
9  & Zero       & \texttt{/dev/zero} \\
10 & PipeRead   & Read end of a pipe \\
11 & PipeWrite  & Write end of a pipe \\
13 & VfsFile    & VFS-backed regular file \\
14 & VfsDir     & VFS-backed directory \\
15 & ProcFs     & Virtual /proc or /sys file \\
16 & TcpSocket  & TCP socket (via net IPC) \\
17 & UdpSocket  & UDP socket (via net IPC) \\
18 & Initrd     & File backed by initrd CPIO \\
22 & EventFd    & eventfd counter \\
23 & TimerFd    & timerfd (TSC-based) \\
25 & MemFd      & memfd (anonymous page-backed) \\
27 & UnixA      & socketpair end A \\
28 & UnixB      & socketpair end B \\
\bottomrule
\end{tabular}
\end{table}


% ------------------------------------------------------------------------------
\chapter{Formal Verification}
\label{ch:formal}
% ------------------------------------------------------------------------------

\section{TLA+ Specifications}

Six TLA+ specifications cover the core kernel subsystems. Each has a
matching TLC model-checking configuration file.

\begin{table}[H]
\centering
\caption{TLA+ specifications and their coverage.}
\begin{tabularx}{\textwidth}{l X}
\toprule
\textbf{Specification} & \textbf{Coverage} \\
\midrule
capabilities        & CDT derivation, revocation, rights attenuation \\
ipc                 & Endpoint send/recv rendezvous, deadlock freedom \\
scheduler           & MLQ priority inversion freedom, work stealing \\
sot\_capabilities    & SOT layer cap derive/revoke with interposition \\
sot\_provenance      & Hash-chain integrity, ring buffer overflow \\
sot\_transactions    & Tier 0/1/2 commit/abort, serialization \\
\bottomrule
\end{tabularx}
\end{table}

All specifications are checked via \texttt{just tlc-mc}, which runs TLC on
each spec and reports per-specification state counts.

\section{sotFS Formal Methods}

The \sotFS{} filesystem uses multiple formal verification approaches:

\begin{description}[leftmargin=1.4em]
  \item[DPO Rewriting.] Each POSIX operation is modeled as a Double Pushout
    rule on the type graph (see Section~\ref{def:dpo-create}).
  \item[Bounded Treewidth.] Proved that $\mathrm{tw}(TG) \leq 6$ without
    hard links, ensuring linear-time graph algorithms.
  \item[Typestate.] Rust's type system enforces valid state transitions at
    compile time (e.g., a transaction handle cannot be used after commit).
  \item[loom.] Concurrency testing for lock-free data structures.
  \item[proptest.] Property-based testing (found 2 bugs: an off-by-one in
    directory iteration and a missing EOF marker).
  \item[cargo-fuzz.] Coverage-guided fuzzing: 300K runs, 0 crashes.
\end{description}

\section{Invariants}

The following invariants are maintained by the system and checked by the
TLA+ specifications:

\begin{invariant}[Capability Monotonicity]
For any capability $c$ derived from $p$:
$\text{rights}(c) \subseteq \text{rights}(p)$.
\end{invariant}

\begin{invariant}[IPC Deadlock Freedom]
No cycle exists in the endpoint wait graph.
\end{invariant}

\begin{invariant}[Transaction Serializability]
The commit order of Tier~2 transactions is equivalent to some serial
execution.
\end{invariant}

\begin{invariant}[Provenance Integrity]
The SHA-256 hash chain in the provenance ring is tamper-evident:
$h_i = \text{SHA256}(h_{i-1} \| \text{event}_i)$.
\end{invariant}

\begin{invariant}[W\^{}X]
For all pages $p$ in any address space:
$\neg(\text{writable}(p) \land \text{executable}(p))$.
\end{invariant}


% ==============================================================================
% APPENDICES
% ==============================================================================
\appendix

% ------------------------------------------------------------------------------
\chapter{Linux ABI Constants}
\label{app:linux-abi}
% ------------------------------------------------------------------------------

The \texttt{sotos-common::linux\_abi} module defines 198
\texttt{SYS\_*} constants for Linux system call numbers. Key structures
defined in this module:

\begin{table}[H]
\centering
\caption{Linux ABI structures in sotos-common.}
\begin{tabularx}{\textwidth}{l c X}
\toprule
\textbf{Struct} & \textbf{Size} & \textbf{Description} \\
\midrule
\texttt{Stat}             & 144\,B & File status (matches Linux x86\_64) \\
\texttt{Termios}          & 44\,B  & Terminal I/O settings \\
\texttt{Winsize}          & 8\,B   & Terminal window size \\
\texttt{KernelSigaction}  & 32\,B  & Signal action (24 + sigsetsize) \\
\texttt{LinuxDirent64}    & var    & Directory entry \\
\bottomrule
\end{tabularx}
\end{table}

Important ABI notes:
\begin{itemize}[leftmargin=1.4em]
  \item \texttt{rt\_sigaction}: The kernel sigaction struct is 32 bytes
    (24 + 8-byte sigset), \emph{not} 152. Writing 152 bytes to
    \texttt{oldact} corrupts stack return addresses.
  \item \texttt{TCGETS} ioctl: reads 36 bytes (not 60).
  \item \texttt{fstat} on kind=14 (VfsDir): returns \texttt{S\_IFDIR}.
\end{itemize}


% ------------------------------------------------------------------------------
\chapter{Project Structure}
\label{app:structure}
% ------------------------------------------------------------------------------

\begin{description}[leftmargin=0em]
  \item[\texttt{kernel/}] The microkernel (Ring~0)
    \begin{description}[leftmargin=1em]
      \item[\texttt{arch/x86\_64/}] CPU-specific: serial, GDT, IDT, LAPIC,
        SMP boot, SYSCALL/SYSRET
      \item[\texttt{mm/}] Physical frame allocator (bitmap), slab allocator
      \item[\texttt{cap/}] Capability table + CDT
      \item[\texttt{ipc/}] Endpoints (sync) + Channels (async)
      \item[\texttt{sched/}] Thread management, MLQ scheduler, Chase-Lev
        deque
    \end{description}
  \item[\texttt{libs/}] Shared libraries
    \begin{description}[leftmargin=1em]
      \item[\texttt{sotos-common/}] Kernel-userspace ABI, syscall numbers,
        linux\_abi module
      \item[\texttt{sotos-pci/}] PCI bus enumeration
      \item[\texttt{sotos-virtio/}] virtio device drivers
      \item[\texttt{sotos-net/}] IP/TCP/UDP/ICMP networking
      \item[\texttt{sotos-nvme/}] NVMe SSD driver
      \item[\texttt{sotos-xhci/}] xHCI USB driver
      \item[\texttt{sotos-objstore/}] Object store + VFS
      \item[\texttt{sotos-ld/}] Dynamic linker
    \end{description}
  \item[\texttt{services/}] Userspace servers
    \begin{description}[leftmargin=1em]
      \item[\texttt{init/}] First userspace process (spawns all others)
      \item[\texttt{vmm/}] Virtual memory manager (paging, CoW)
      \item[\texttt{kbd/}] Keyboard driver service
      \item[\texttt{net/}] Network service
      \item[\texttt{nvme/}] NVMe driver service
      \item[\texttt{xhci/}] xHCI USB driver service
    \end{description}
  \item[\texttt{personality/}] OS personality layers
    \begin{description}[leftmargin=1em]
      \item[\texttt{bsd/}] NetBSD rump kernel integration
      \item[\texttt{linux/}] Linux compatibility (LUCAS, LKL)
      \item[\texttt{common/deception/}] Deception engine
    \end{description}
  \item[\texttt{boot/}] Bootloader configuration (\texttt{limine.conf})
  \item[\texttt{scripts/}] Build tools (\texttt{mkimage.py}, test scripts)
  \item[\texttt{docs/}] Documentation and specifications
  \item[\texttt{specs/}] TLA+ formal specifications
\end{description}


% ------------------------------------------------------------------------------
\chapter{VfsOp to DPO Rule Mapping}
\label{app:vfsop}
% ------------------------------------------------------------------------------

Table~\ref{tab:vfsop-dpo} provides the complete mapping from VFS operations
to their DPO rewriting rules, including the graph elements affected.

\begin{longtable}{l l l l l}
\caption{Complete VfsOp to DPO rule mapping.} \label{tab:vfsop-dpo} \\
\toprule
\textbf{VfsOp} & \textbf{L (match)} & \textbf{K (keep)} & \textbf{R (produce)} & \textbf{Gluing} \\
\midrule
\endfirsthead
\toprule
\textbf{VfsOp} & \textbf{L (match)} & \textbf{K (keep)} & \textbf{R (produce)} & \textbf{Gluing} \\
\midrule
\endhead
\bottomrule
\endlastfoot
CREATE & Dir $d$ & Dir $d$ & $d$, Inode $i$, $d \to i$ & $d$ identity \\
MKDIR  & Dir $d$ & Dir $d$ & $d$, Dir $d'$, $d \to d'$, $d' \to d$ & $d$ identity \\
RMDIR  & Dir $d$, Dir $d'$, $d \to d'$ & Dir $d$ & Dir $d$ & empty $d'$ \\
LINK   & Dir $d$, Inode $i$ & $d$, $i$ & $d$, $i$, $d \to i$ & $d$, $i$ id \\
UNLINK & Dir $d$, $d \to i$ & Dir $d$ & Dir $d$ & nlink check \\
RENAME & Dir $s$, $s \to i$, Dir $t$ & $s$, $i$, $t$ & $s$, $i$, $t$, $t \to i$ & $s$, $t$ id \\
WRITE  & Inode $i$ & Inode $i$ & $i$, Ver $v$, Blk $b$, edges & $i$ identity \\
\end{longtable}


% ------------------------------------------------------------------------------
\chapter{Configuration Reference}
\label{app:config}
% ------------------------------------------------------------------------------

\section{QEMU Flags}

\begin{longtable}{l p{8cm}}
\caption{Common QEMU flags used with sotBSD.} \\
\toprule
\textbf{Flag} & \textbf{Purpose} \\
\midrule
\endfirsthead
\endhead
\bottomrule
\endlastfoot
\texttt{-serial stdio}         & Redirect COM1 to terminal \\
\texttt{-display none}         & Headless (serial-only) \\
\texttt{-no-reboot}            & Halt on triple fault \\
\texttt{-m 2048M}              & 2\,GiB RAM \\
\texttt{-smp N}                & N CPU cores \\
\texttt{-accel whpx}           & Windows Hypervisor (fast) \\
\texttt{-accel tcg}            & Software emulation (slow) \\
\texttt{-s}                    & GDB server on port 1234 \\
\texttt{-device virtio-net-pci}& virtio network adapter \\
\texttt{-device nvme}          & NVMe SSD controller \\
\texttt{-device qemu-xhci}     & xHCI USB 3.0 controller \\
\texttt{-netdev user,id=n0}    & SLIRP user-mode networking \\
\end{longtable}

\section{Key Compile Flags}

\begin{description}[leftmargin=1.4em]
  \item[\texttt{-C relocation-model=static}] Required: Limine rejects
    \texttt{ET\_DYN}.
  \item[\texttt{-Z build-std=core,alloc}] Build \texttt{core} and
    \texttt{alloc} from source for the bare-metal target.
  \item[\texttt{-C code-model=kernel}] Place code in the upper 2\,GiB
    (optional, for HHDM compatibility).
  \item[\texttt{CARGO\_ENCODED\_RUSTFLAGS}] Override parent
    \texttt{.cargo/config.toml} for userspace crate builds.
\end{description}


% ==============================================================================
% BACK MATTER
% ==============================================================================

\chapter*{Glossary}
\addcontentsline{toc}{chapter}{Glossary}

\begin{description}[leftmargin=0em]
  \item[AP] Application Processor (non-boot CPU in SMP).
  \item[BSP] Bootstrap Processor (the CPU that runs firmware).
  \item[CDT] Capability Derivation Tree.
  \item[CoW] Copy-on-Write.
  \item[DPO] Double Pushout (graph rewriting framework).
  \item[FMA] Fault Management Architecture.
  \item[HHDM] Higher Half Direct Map.
  \item[IPC] Inter-Process Communication.
  \item[IST] Interrupt Stack Table.
  \item[LAPIC] Local Advanced Programmable Interrupt Controller.
  \item[LUCAS] Linux User-mode Compatibility And Syscall translation.
  \item[MLQ] Multi-Level Queue (scheduler).
  \item[PRP] Physical Region Page (NVMe scatter-gather).
  \item[SMF] Service Management Facility.
  \item[SO] Secure Object.
  \item[SOT] Secure Object Transaction.
  \item[SPSC] Single-Producer Single-Consumer (ring buffer).
  \item[TLC] TLA+ model checker.
  \item[UC] Uncacheable (memory type for MMIO).
  \item[vDSO] virtual Dynamic Shared Object.
  \item[VNIC] Virtual Network Interface Card.
  \item[VFS] Virtual File System.
  \item[VMM] Virtual Memory Manager.
  \item[W\^{}X] Write XOR Execute.
\end{description}

\end{document}
